\section{The closed arithmetic source and its spectral scale}
\label{sec:closed-source}
Finite-dimensional ground-state geometry does not supply the full
arithmetic input space. We now construct a closed infinite-dimensional
preparation for the gamma kinetic energy, retaining all remaining
arithmetic terms explicitly.

\subsection{Local fields, maximal domain and adjoint}
Let $H=-\partial_x^2$ on $L^2(\R)$, $a_k=2k+1/2$, and
$R_k=(a_k^2+H)^{-1}$. For a prescribed source $F$, minimize
\[
\mathcal E_k(F,u)=\frac2{a_k}
 \left(\norm{F-u}^2+a_k^{-2}\norm{u'}^2\right)
\quad\text{over }u\in H^1(\R).
\]
The minimizer is $u_k=a_k^2R_kF$. Define the output space and residuals
\begin{equation}\label{source:definition}
\mathscr Y_\gamma=\bigoplus_{k\ge0}(L^2(\R)\oplus L^2(\R)),
\qquad
(\mathcal A_\gamma F)_k=
\begin{pmatrix}
\sqrt{2/a_k}\,H R_kF\\
\sqrt{2a_k}\,\partial_xR_kF
\end{pmatrix}.
\end{equation}
The Fourier column $v_k(\tau)$ obtained from this expression has
$v_k^\dagger v_k=(2/a_k)\tau^2/(a_k^2+\tau^2)$. Hence
\[
\norm{\mathcal A_\gamma F}^2
=\frac1{2\pi}\int_\R B(\tau^2)|\widehat F(\tau)|^2\,d\tau=K[F].
\]
The digamma partial fraction and asymptotic formulas \cite{DLMF} give
\begin{equation}\label{source:log}
B(\tau^2)=\re\psi(1/4+i\tau/2)-\psi(1/4)
=\log|\tau|-\log2-\psi(1/4)+O(\tau^{-2}).
\end{equation}
The positive measure
$\sum_k(2/a_k)\delta_{a_k^2}$ also places this explicit tower in the
complete-Bernstein extension framework \cite{KwasnickiMucha}.

For $X_L=L^2(I_L)$ and zero extension $E_L$, write
$A_L=\mathcal A_\gamma E_L$ and
\begin{equation}\label{source:domain}
\mathcal D_{\log,L}=
\left\{f\in X_L:
 \int_\R B(\tau^2)|\widehat{E_Lf}(\tau)|^2\,d\tau<\infty\right\}.
\end{equation}
This is the maximal source domain. It is equivalent to the
Fourier weight $1+\log(2+|\tau|)$ after adding the $L^2$ norm.

\begin{proposition}\label{source:closed}
The source $A_L$ is closed and densely defined, with domain
$\mathcal D_{\log,L}$ and core $C_c^\infty(I_L)$.
\end{proposition}
\begin{proof}
Each component of \eqref{source:definition} is bounded. If
$f_n\to f$ and $A_Lf_n\to y$ in the two Hilbert spaces, componentwise
convergence identifies every component of $y$ with that of $A_Lf$.
Their squared sum is finite, proving closedness.
For the core assertion, first contract the supported function using
$D_rf(x)=r^{-1/2}f(x/r)$, $r<1$. These maps are uniformly bounded
near $r=1$ in the logarithmic Fourier norm and converge strongly
to the identity there, by density and the weight comparison under
dilation. The contracted support has a positive gap from the
endpoints. Mollification inside that gap then converges in the
same norm.
\end{proof}

The maximal adjoint must be taken after interval restriction.
For $\Psi\in\mathscr Y_\gamma$, define the distribution $t_\Psi$ by
\[
\widehat t_\Psi(\tau)=\sum_k v_k(\tau)^\dagger\widehat\Psi_k(\tau).
\]
Cauchy--Schwarz gives
$|\widehat t_\Psi|^2/(1+B)\le\sum_k|\widehat\Psi_k|^2$, so the sum
defines a distribution. The global adjoint requires $t_\Psi\in L^2(\R)$,
whereas
\begin{equation}\label{source:adjoint}
\Dom A_L^*=\{\Psi:t_\Psi|_{I_L}\in L^2(I_L)\},
\qquad A_L^*\Psi=t_\Psi|_{I_L}.
\end{equation}
In particular the interval adjoint domain can be larger. The form
operator $T_L=A_L^*A_L$ has domain
\begin{equation}\label{source:operator-domain}
\Dom T_L=
\{f\in\mathcal D_{\log,L}:
 (B(H)E_Lf)|_{I_L}\in L^2(I_L)\},
\end{equation}
where the multiplier is first understood distributionally.
The block charge with entries $A_L,A_L^*$ on
$\Dom A_L\oplus\Dom A_L^*$ is self-adjoint, and its square has
diagonal entries $T_L,A_LA_L^*$. This construction concerns the
kinetic term.

\subsection{Two controlled singular limits}
Let $A_{L,N}$ retain $0\le k<N$. Its exact operator norm is
\begin{equation}\label{source:tower-norm}
\norm{A_{L,N}}^2
=\sum_{k<N}\frac2{a_k}
=\psi(N+1/4)-\psi(1/4)
=\log N-\psi(1/4)+O(N^{-1}).
\end{equation}
The upper bound follows from the multiplier; modulating any fixed
nonzero compactly supported input to arbitrarily high frequency
attains its supremum, including after compression. Increasing
towers converge in output norm precisely on
$\mathcal D_{\log,L}$; outside it their squared norms diverge.
If $F=E_Lf\in H^1(\R)$, the tail obeys
\[
\norm{(A_L-A_{L,N})f}^2
\le\left(\frac2{a_N^3}+\frac1{2a_N^2}\right)\norm{F'}^2,
\]
where $A_{L,N}$ is embedded by zero components.
Monotone convergence of closed forms gives strong-resolvent
convergence of the associated operators \cite{SimonForms}.

The heat regulator is
$A_{L,\epsilon}=\mathcal A_\gamma e^{-\epsilon H/2}E_L$.
Its point-source columns are the inverse Fourier transforms of
$v_k(\tau)e^{-\epsilon\tau^2/2}e^{-i\tau x}$, and
\begin{align}
\norm{A_{L,\epsilon}}_{\mathrm{HS}}^2
&=\frac L{2\pi}\int_\R B(\tau^2)e^{-\epsilon\tau^2}\,d\tau
\label{source:heat}\\
&=\frac{L}{2\sqrt{\pi\epsilon}}
\left(\frac12\log\frac1\epsilon+
\frac{\gamma_{\!E}}2+\frac\pi2+\log2\right)+O_L(1).
\nonumber
\end{align}
To obtain the constant, integrate the logarithm in
\eqref{source:log}, using
$\psi(1/2)=-\gamma_{\!E}-2\log2$ and
$\psi(1/4)=-\gamma_{\!E}-\pi/2-3\log2$.
The remainder is integrable after subtracting the high-frequency
asymptotic with a cutoff. The source and form limits again have
exactly domain $\mathcal D_{\log,L}$, with strong-resolvent
convergence. Thus a regularized family of normalizable point caps
has a necessarily singular size limit.

\subsection{The complete signed form}
Let $c(x)=\cosh(x/2)$, $s(x)=\sinh(x/2)$ on $I_L$, and
$T_d=E_L^*U_dE_L$. Define the bounded self-adjoint remainder
\begin{equation}\label{source:remainder}
R_L=w_0I+2cc^*-2ss^*
 -\sum_{m\log p<L}(\log p)p^{-m/2}(T_{m\log p}+T_{m\log p}^*).
\end{equation}
Here $cc^*$ denotes $f\mapsto c\ip c f$; the squared integrals
in the original convention have the same quadratic values.
A useful coarse bound is
\[
\norm{R_L}\le |w_0|+L+2\sinh(L/2)
                  +2\sum_{m\log p<L}(\log p)p^{-m/2}.
\]
The complete operator is
$W_L=T_L+R_L$, with form domain $\mathcal D_{\log,L}$ and
operator domain $\Dom T_L$. Adding this same $R_L$ to either
regulator preserves the complete signed expression and the
strong-resolvent limit. These operators are uniformly bounded
below; that fact is not their nonnegativity.
If $I_L\subset I_{L_*}$, zero extension preserves the complete
quadratic form, since all newly available translations with no
overlap vanish.

The embedding $\mathcal D_{\log,L}\hookrightarrow X_L$ is compact.
Indeed, the weighted norm uniformly controls the high-frequency
tail, while restricting a fixed finite Fourier band to a finite
interval is Hilbert--Schmidt. Thus $T_L$ and $W_L$ have compact
resolvent. There are only finitely many negative eigenvalues of
$W_L$ at any fixed $L$, without a claim about their absence.

\subsection{Why bounded and finite-rank sources do not suffice}
For $\varphi\in C_c^\infty(I_L)$ and
$f_N(x)=e^{iNx}\varphi(x)$, expansion of the multiplier gives
\begin{equation}\label{source:modulation}
\begin{split}
Q_L[f_N]={}&\norm\varphi^2(\log N-\log(2\pi))\\
&-2\sum_{m\log p<L}(\log p)p^{-m/2}
\re\!\left(e^{-iNm\log p}
\ip{E_L\varphi}{U_{m\log p}E_L\varphi}\right)+o(1).
\end{split}
\end{equation}
The pole amplitudes decay faster than any power. In particular
$Q_L[f_N]=\norm\varphi^2\log N+O(1)$. No bounded source can
represent this form. A square-integrable family of point caps
$b_x$ would define a Hilbert--Schmidt source and is also excluded.
Evolution confined to a finite vacuum space, on which $H=0$,
does not supply these inputs.

Nor can an unbounded finite-range source evade the obstruction by
a domain restriction. Every densely defined closable operator
with finite-dimensional range is bounded on its domain: otherwise
choose inputs tending to zero whose outputs have norm one, then
take a convergent output subsequence, contradicting closability.
Its closure is therefore bounded and everywhere defined.
If a proposed norm identity holds on the smooth core, the closed
arithmetic form makes that source closable: a sequence of inputs
tending to zero with Cauchy outputs is Cauchy in the shifted form
norm, so the closed-form limit forces the output limit to vanish.
Equation \eqref{source:modulation} excludes the resulting bounded
operator.

\subsection{A resolvent pairing as a general control}
For an independently self-adjoint physical Hamiltonian $\mathcal H$
and bounded finite-channel coupling $J$, put
$M(z)=J^*(\mathcal H-z)^{-1}J$ for $\im z>0$ and
$\Gamma(z)=(\mathcal H-\bar z)^{-1}J$. The resolvent identity gives
\begin{equation}\label{source:nevanlinna}
\frac{M(z)-M(w)^*}{z-\bar w}=\Gamma(z)^*\Gamma(w).
\end{equation}
This is a positive Gram kernel. Its Cayley transform
$S(z)=(M(z)-iI)(M(z)+iI)^{-1}$ satisfies
\[
\frac{I-S(z)S(w)^*}{-i(z-\bar w)}
=2(M(z)+iI)^{-1}
 \frac{M(z)-M(w)^*}{z-\bar w}(M(w)^*-iI)^{-1}.
\]
In a general operator Herglotz representation
\[
M(z)=A+Dz+\int_\R
 \left(\frac1{t-z}-\frac{t}{1+t^2}\right)d\Sigma(t),
\quad D\ge0,
\]
the kernel sees $D$ and the positive measure but loses the constant
$A$ \cite{Herglotz}. This illustrates both the strength and the
normalization limit of a resolvent Gram law. A regular bounded
coupling does not on its own identify the complete arithmetic form
or its singular source limit.

\subsection{Neumann--Dirichlet comparison}
Let $H_{\rm N},H_{\rm D}$ be the ordinary Neumann and Dirichlet
interval Laplacians.
\begin{theorem}\label{source:bracket}
In the sense of closed quadratic forms,
\[
B(H_{\rm N})\le T_L\le B(H_{\rm D}),
\]
with
$\Dom B(H_{\rm D})^{1/2}\subset\mathcal D_{\log,L}
\subset\Dom B(H_{\rm N})^{1/2}$. Consequently, indexing from one,
\begin{equation}\label{source:eigenvalues}
B(((n-1)\pi/L)^2)\le\lambda_n(T_L)\le B((n\pi/L)^2),
\qquad
\lambda_n(T_L)=\log\frac{\pi n}{2L}-\psi(1/4)+O_L(n^{-1}).
\end{equation}
\end{theorem}
\begin{proof}
For each mass, restricting a whole-line auxiliary field to $I_L$
drops nonnegative exterior energy; extending a Dirichlet field
by zero gives an admissible whole-line competitor. Minimizing
the three energies therefore orders them as Neumann, whole-line,
Dirichlet. The outer minima are the forms of
$(2/a)H_{\rm N,D}(a^2+H_{\rm N,D})^{-1}$.
Sum finite towers and pass to increasing form limits. Min--max
and the known interval Laplacian eigenvalues give the bounds;
\eqref{source:log} gives the asymptotic.
\end{proof}
This comparison concerns three different preparations. It does
not impose Neumann or Dirichlet boundary conditions on the actual
whole-line gamma fields.
